\documentclass[aspectratio=169]{beamer}
\usepackage[utf8]{inputenc}\usepackage{amsmath}\usepackage{tcolorbox}
\definecolor{UNIBBlue}{RGB}{0, 51, 102}\setbeamercolor{frametitle}{bg=UNIBBlue,fg=white}
\title{Optimisasi untuk Teknik Elektro}\subtitle{Minggu 9: Mixed Integer Linear Programming (MILP)}
\date{Semester Ganjil 2026}\begin{document}\begin{frame}\titlepage\end{frame}
\begin{frame}{Variabel Diskrit dalam Keputusan}
Dalam teknik elektro, banyak keputusan bersifat \textit{diskrit} atau logika \textit{Ya/Tidak}.
\begin{itemize}
\item Apakah pembangkit dihidupkan (1) atau dimatikan (0)? (Unit Commitment).
\item Jumlah kabel yang harus dibentangkan (bilangan bulat).
\end{itemize}
MILP memungkinkan kita memodelkan logika sekuensial dan pemilihan ini secara eksak, sering diselesaikan dengan metode \textit{Branch-and-Bound}.
\end{frame}\end{document}